A segmentation-and-refit framework exhibits a related separation between
pooled and temporal gains (Table~\ref{tab:external-framework}). On Study C's
30 paired repeats at each duration and the $300/150$ grid, penalised PELT
partitions the raw window features into contiguous segments. A VAR-Lasso
model is then fitted independently on each non-overlapping segment range.
This comparison uses our graph-estimation adapter to the external change-point
algorithm; it estimates segment count and uses neither recurrent K-means
grouping nor the acceptance gate.

Relative to raw windows, PELT plus VAR increases pooled AUROC by $+0.013$
$[0.006,0.020]$ at $D=300$ and $+0.008$ $[0.001,0.014]$ at
$D=600$. Within-edge temporal AUROC differences are $+0.005$
$[-0.020,0.027]$ and $-0.006$ $[-0.020,0.008]$, leaving the temporal
comparison inconclusive at both durations. At $D=600$, positive-change AP
gives $-0.125$ $[-0.205,-0.044]$ while standard AP gives $-0.011$
$[-0.067,0.048]$. Thus both the evaluation target and zero-score convention
affect the comparison outside recurrent clustering.

All 60 cases and 644 segment refits converge. The fixed configuration yields
9--11 segments on 23 windows, reaching the 11-segment minimum-length cap in
45 cases. Supplement Section~S6 provides the settings, counts and paired
results, alongside the distinct auxiliary seven-window segmentation study.
